\documentclass[11pt]{article}
\usepackage{amsmath,amssymb,amsthm,fullpage,hyperref}
\newtheorem{proposition}{Proposition}
\title{Audit of a Discrete-Sigmoid Degree Bound}
\author{}
\date{November 24, 2026}
\begin{document}
\maketitle

\paragraph{Status.}
This is an independent written audit of a potentially useful literature
lead. It proves that one polynomial-degree step in the printed argument
cannot hold as stated. It does not disprove the printed VC-dimension
theorem, and it proves neither accepted HeadComplexity target. Nothing
here changes the model or is imported as a Lean theorem.

\section{The lead and its limitation}
Anthony and Bartlett, \emph{Neural Network Learning: Theoretical
Foundations}, Theorem 8.11, printed page 122, gives
$2W\log_2(60ND)$ as an upper bound for a two-layer standard-sigmoid
network with $W$ parameters, integer inputs in $[-D,D]$, and first-layer
fan-in at most $N$. Its proof on printed page 123 transforms each hidden
unit to a rational function of degree at most $ND$ in the numerator and
$2ND+1$ in the denominator, then asserts a degree bound $3ND+2$ for
the polynomial whose sign is the output. The preceding Theorem 8.3
uses total degree, explicitly defined on printed page 108 as the sum
of monomial exponents. Thus a coordinatewise-degree convention does
not explain the missing dependence on the number $k$ of hidden units.

Straight denominator clearing gives the valid bound
$k(2ND+1)+1$: each output summand has its own numerator and all the
other denominator factors. This returns an order
$W\log(kND)$ VC estimate. The abstract of the primary 1996 paper by
Bartlett and Williamson reports order $W\log(WD)$, consistent with
this dependence. No erratum resolving the discrepancy was found in
the author's book webpage or targeted searches during this audit.

\section{The degree cannot be removed at this proof step}
Even at the single allowed input $x=0$, a sigmoid network contains
the parameter family
\[
 Z(a,u,b)=b+\sum_{j=1}^k\frac{u_j}{1+a_j},
 \qquad a_j>0,\quad u_j,b\in\mathbb R.
\]
Here $a_j$ is the exponential of the $j$th hidden bias. Input weights
may be fixed at zero; all output weights remain unrestricted.

\begin{proposition}
If a nonzero real polynomial $Q(a,u,b)$ has the same strict sign as
$Z$ wherever $Z\ne0$ on this parameter domain, then
$\deg Q\ge k+1$.
\end{proposition}
\begin{proof}
Set $L_j=1+a_j$, $D=\prod_jL_j$, and
\[
 P=bD+\sum_{j=1}^ku_j\prod_{i\ne j}L_i.
\]
Its total degree is $k+1$. Regard $P$ as a polynomial in $b$ over
the unique factorization domain $R=\mathbb R[a_1,\ldots,a_k,
u_1,\ldots,u_k]$. Its leading coefficient $D$ is relatively prime
to its constant coefficient. Indeed any irreducible factor of $D$
is some $L_j$, and reduction of the constant coefficient modulo
$L_j$ is $u_j\prod_{i\ne j}L_i$, which is nonzero. Therefore $P$
is primitive. It is linear and hence irreducible over the fraction
field of $R$; Gauss's lemma makes it irreducible in $R[b]$.

For every positive $a$ and every real $u$, variation of $b$ crosses
the zero of $Z$ with positive slope. Continuity and the assumed
strict signs force
\[
 Q\left(a,u,-\sum_j u_j/L_j\right)=0.
\]
After clearing a power of $D$, this is a polynomial identity: a
polynomial vanishing on the open set
$(0,\infty)^k\times\mathbb R^k$ vanishes identically. Consequently
$P$ divides $Q$ over the fraction field of $R$ and, by primitivity,
in $R[b]$. Total degrees add for nonzero polynomial products, giving
the claim.
\end{proof}

For $N=D=1$ and $k\ge5$, this contradicts the proof's claimed
degree bound $3ND+2=5$ in precisely its proposed transformed
parameters. This is a failure of that proof step, not a
counterexample to its VC-dimension conclusion. At the single input
used above there is of course no large shattered set. A different
sign-counting argument could conceivably prove the stated bound,
but none is supplied by the cited polynomial-degree calculation.

\section{Consequence for the present research}
The bound cannot currently be used as a shortcut to an improved
unrestricted count for sums of admissible reciprocals. Positive
denominators do not by themselves justify dropping the number of
summands from the total degree. The missing task remains a direct
sign-pattern argument exploiting additional structure, or a
different construction for the upper-bound target.

\begin{thebibliography}{9}
\bibitem{AB}
M. Anthony and P. L. Bartlett,
\emph{Neural Network Learning: Theoretical Foundations},
Cambridge University Press, 1999, Chapter 8.
\href{https://www.stat.berkeley.edu/~bartlett/nnl/index.html}{Author's book webpage}.
\href{https://openeclass.panteion.gr/modules/document/file.php/PMS152/LEARNING/Anthony\%20Martin\%2C\%20Bartlett\%20Peter\%20L.\%20\%282009\%29\%20--\%20Neural\%20Network\%20Learning_\%20Theoretical\%20Foundations.pdf}{Consulted scan}.
\bibitem{BW}
P. L. Bartlett and R. C. Williamson,
The VC dimension and pseudodimension of two-layer neural networks
with discrete inputs, \emph{Neural Computation} 8(3), 625--628, 1996.
\href{https://doi.org/10.1162/neco.1996.8.3.625}{doi:10.1162/neco.1996.8.3.625}.
\href{https://ieeexplore.ieee.org/document/6796931/}{Publisher abstract}.
\end{thebibliography}
\end{document}
